\section{模型假设}

为保证模型能够在题目给定数据范围内求解，并使种植方案具有明确的农业含义，本文作如下假设。

\begin{enumerate}
  \item 假设 2023 年各作物在各季次的实际总产量可作为 2024--2030 年的基准预期销售量。即在未考虑未来市场变化时，某作物某季次的预期销售量等于 2023 年该作物该季次在全村范围内的实际产量。该假设能够充分利用附件2中已有的真实种植记录，并避免额外引入难以验证的外部需求数据。
  \item 假设亩产量、种植成本和销售价格由作物、地块类型和季次共同决定。附件2中销售价格以区间形式给出，本文取区间中点作为基准销售单价；智慧大棚相对于普通生产条件的价格差异，按题目统计表中的地块类型口径自动体现在价格参数中，不再额外叠加修正。
  \item 假设智慧大棚第一季蔬菜的亩产量和种植成本可由普通大棚第一季同作物数据平补。附件2中缺少智慧大棚第一季对应统计数据，而智慧大棚第一季实际可种植蔬菜，若直接删除该类组合会造成可行集合不完整。由于普通大棚与智慧大棚在第一季蔬菜种植上具有较强可比性，故采用普通大棚第一季数据作为平补依据。
  \item 假设种植面积和是否种植为规划期内的核心决策。种植面积用于刻画资源分配，0-1种植指示变量用于表达最小种植比例、连作禁止、豆类轮作和作物分散度等离散约束。问题二和问题三中，种植拓扑作为第一阶段决策，各随机情景下的销售分配作为第二阶段决策。
  \item 假设未来不确定性主要来自销售量、种植成本、销售价格、亩产波动和露天灾害。小麦、玉米需求按题目给定增长区间变化，其他作物需求在给定范围内波动；种植成本逐年上升；蔬菜价格呈增长趋势，食用菌价格呈下降趋势；亩产量存在随机波动。露天耕地受自然灾害影响，普通大棚和智慧大棚因具有设施保护，不纳入灾害减产范围。
  \item 假设作物之间的替代关系可由市场指标的负相关性近似刻画，互补关系主要体现为豆类前茬对后茬作物成本的降低。本文利用 Spearman 相关系数和聚类结果识别同类作物中的替代关系，并将豆类轮作带来的土壤改良作用转化为后茬成本折减系数。
\end{enumerate}

\section{符号说明}

本文主要符号如表\ref{tab:symbols}所示。其余局部变量将在相应模型建立过程中说明。

\begin{longtable}{p{0.16\textwidth}p{0.56\textwidth}p{0.18\textwidth}}
\caption{主要符号说明}\label{tab:symbols}\\
\toprule
符号 & 含义 & 单位或维度 \\
\midrule
\endfirsthead
\toprule
符号 & 含义 & 单位或维度 \\
\midrule
\endhead
$T$ & 规划年份集合，$T=\{2024,\ldots,2030\}$ & 7 年 \\
$I$ & 地块集合 & 54 个地块 \\
$J$ & 作物集合 & 41 种作物 \\
$S$ & 季次集合，$S=\{1,2\}$ & 2 个季次 \\
$\Omega$ & 可行种植三元组集合 $(i,j,s)$ & 地块--作物--季次 \\
$i,j,s,t$ & 地块、作物、季次、年份编号 & 索引 \\
$k(i)$ & 地块 $i$ 的地块类型 & 6 类地块 \\
$A_i$ & 地块 $i$ 的面积 & 亩 \\
$Y_{i,j,s}$ & 作物 $j$ 在地块 $i$、季次 $s$ 的亩产量 & 斤/亩 \\
$C_{i,j,s}$ & 作物 $j$ 在地块 $i$、季次 $s$ 的种植成本 & 元/亩 \\
$P_{i,j,s}$ & 作物 $j$ 在地块 $i$、季次 $s$ 的销售单价 & 元/斤 \\
$D_{j,s}$ & 作物 $j$ 在季次 $s$ 的基准预期销售量 & 斤 \\
$x_{i,j,s,t}$ & 年份 $t$ 在地块 $i$、季次 $s$ 种植作物 $j$ 的面积 & 亩 \\
$u_{i,j,s,t}$ & 年份 $t$ 在地块 $i$、季次 $s$ 是否种植作物 $j$ & 0-1 变量 \\
$q^n_{i,j,s,t}$ & 正常价格销售的产量 & 斤 \\
$q^d_{i,j,s,t}$ & 折价销售的产量 & 斤 \\
$\gamma$ & 超产部分销售折扣系数 & 0 或 0.5 \\
$\delta$ & 最小种植比例 & 0.10 \\
$N_j^{\max}$ & 作物 $j$ 可分散种植的地块数上限 & 个 \\
$J_{\rm bean}$ & 豆类作物集合 & 8 种作物 \\
$W$ & 问题一确定性总利润 & 元 \\
$\omega$ & 随机情景编号 & $\omega=1,\ldots,K$ \\
$K$ & 随机情景数量 & 基准取 50 \\
$D_{j,s,t}(\omega)$ & 情景 $\omega$ 下作物 $j$ 在年份 $t$、季次 $s$ 的销售量上限 & 斤 \\
$Y_{i,j,s,t}(\omega)$ & 情景 $\omega$ 下的亩产量 & 斤/亩 \\
$C_{i,j,s,t}(\omega)$ & 情景 $\omega$ 下的种植成本 & 元/亩 \\
$P_{i,j,s,t}(\omega)$ & 情景 $\omega$ 下的销售单价 & 元/斤 \\
$\Pi_t(\omega)$ & 年份 $t$、情景 $\omega$ 下的利润 & 元 \\
$\alpha$ & CVaR 置信水平 & 0.95 \\
$\lambda$ & 风险厌恶系数 & 0.10 \\
$\eta_t$ & CVaR 线性化中的分位点辅助变量 & 元 \\
$v_t(\omega)$ & CVaR 线性化中的尾部损失辅助变量 & 元 \\
$\rho_{j,m}$ & 作物 $j$ 与作物 $m$ 的 Spearman 相关系数 & 无量纲 \\
$\mathrm{Sub}(j)$ & 作物 $j$ 的替代作物集合 & 集合 \\
$\theta_{m,j}$ & 作物 $m$ 向作物 $j$ 的需求转移系数 & 无量纲 \\
$\phi$ & 豆类前茬带来的后茬成本折减系数 & 0.90 \\
\bottomrule
\end{longtable}
